\chapter{Complexity — NP}
%%% generated by the blueprint pipeline from TCSlib.Complexity.ClassNP.NP
%%%%%%%%%%%%%%%%%%%%%%%%%%%%%%%%%%%%%%%%%%%%%%%%%%%%%%%%%%%%%%%%%%%%%%%%%%%%%%%
\section{Overview}

This module defines the class $\mathsf{NP}$ following Arora--Barak
Definition~2.1: a language over the binary alphabet is in $\mathsf{NP}$ when
membership is witnessed by certificates of length exactly $C \cdot (|x|+1)^c$
whose concatenation with the input is accepted by a verifier language in
$\mathsf{P}$. It proves $\mathsf{P} \subseteq \mathsf{NP}$ and shows
(Arora--Barak Exercise~2.1) that bounded-length certificates, paired with the
input through a self-delimiting encoding, define the same class; the supporting
material covers marker padding and stripping of certificates, the bounded search
that recovers the input/certificate boundary, and polynomial-time deciders for
the two converted verifier languages.

%%%%%%%%%%%%%%%%%%%%%%%%%%%%%%%%%%%%%%%%%%%%%%%%%%%%%%%%%%%%%%%%%%%%%%%%%%%%%%%
\section{Declarations}

\begin{definition}[The class $\mathsf{NP}$]
\lean{Complexity.NP}
\statementsource{BGS75}{
  pdf-2459e3840724-p002-b002,
  pdf-2459e3840724-p002-b003,
  pdf-2459e3840724-p010-b004,
  pdf-2459e3840724-p010-b005
}
\label{Complexity.NP}
\leanok
\uses{Complexity.P}
A language $L \subseteq \{0,1\}^{*}$ belongs to the class $\mathsf{NP}$ if there
exist natural numbers $C$ and $c$ and a verifier language $V \in \mathsf{P}$ such
that for every binary string $x$, $x \in L$ holds if and only if there is a
certificate $u \in \{0,1\}^{*}$ of length exactly $|u| = C \cdot (|x|+1)^c$ whose
concatenation $xu$ with $x$ lies in $V$.
\end{definition}

\begin{theorem}[Polynomial time is contained in nondeterministic polynomial time]
\lean{Complexity.P_subset_NP}
\label{Complexity.P_subset_NP}
\leanok
\uses{Complexity.NP, Complexity.P}
\difficulty{3}
Every language in $\mathsf{P}$ belongs to $\mathsf{NP}$, where membership in
$\mathsf{NP}$ means that there are natural numbers $C,c$ and a verifier language
$V\in\mathsf{P}$ such that, for each binary string $x$, membership of $x$ is equivalent
to the existence of a binary certificate $u$ of length exactly $C(|x|+1)^c$ for which
the concatenation $xu$ belongs to $V$.
\end{theorem}

\begin{definition}[Stripping a marked certificate]
\lean{Complexity.stripCertificate}
\label{Complexity.stripCertificate}
\leanok
The partial map $\mathrm{strip}$ on binary strings that deletes the last $1$
together with the run of $0$s following it. It is defined recursively:
$\mathrm{strip}(\varepsilon)$ is undefined, and for a bit $b$ and a string $v$,
$\mathrm{strip}(bv) = b\,\mathrm{strip}(v)$ when $\mathrm{strip}(v)$ is defined,
while if $\mathrm{strip}(v)$ is undefined then $\mathrm{strip}(1v) = \varepsilon$
and $\mathrm{strip}(0v)$ is undefined. Thus $\mathrm{strip}(v)$ is the prefix of
$v$ preceding its last $1$, and is undefined when $v$ contains no $1$.
\end{definition}

\begin{lemma}[Stripping an all-zero string]
\lean{Complexity.stripCertificate_false}
\label{Complexity.stripCertificate_false}
\leanok
\uses{Complexity.stripCertificate}
\difficulty{3}
For every $k\in\mathbb{N}$, stripping the binary string consisting of $k$ zeros is
undefined.
\end{lemma}

\begin{lemma}[Stripping a padded string]
\lean{Complexity.stripCertificate_pad}
\label{Complexity.stripCertificate_pad}
\leanok
\uses{Complexity.stripCertificate, Complexity.stripCertificate_false}
\difficulty{3}
For every binary string $u$ and every $k\in\mathbb{N}$, stripping the string formed by
appending a $1$ and then $k$ zeros to $u$ returns $u$.
\end{lemma}

\begin{lemma}[Characterisation of certificate stripping]
\lean{Complexity.stripCertificate_spec}
\label{Complexity.stripCertificate_spec}
\leanok
\uses{Complexity.stripCertificate, Complexity.stripCertificate_false, Complexity.stripCertificate_pad}
\difficulty{5}
For every binary string $v$, stripping $v$ is undefined if and only if $v$ consists
entirely of zeros. For every binary string $u$, stripping $v$ returns $u$ if and only if
$v$ is $u$ followed by a $1$ and then some finite number of zeros.
\end{lemma}

\begin{lemma}[Strict increase of padded length]
\lean{Complexity.certificateTotal_strictMono}
\label{Complexity.certificateTotal_strictMono}
\leanok
\difficulty{3}
For all $C,c\in\mathbb{N}$, the function $n\mapsto n+(C+1)(n+1)^c$ on $\mathbb{N}$ is
strictly increasing.
\end{lemma}

\begin{lemma}[Room for a marker bit]
\lean{Complexity.certificate_room}
\label{Complexity.certificate_room}
\leanok
\difficulty{3}
For all $C,c,n\in\mathbb{N}$, $C(n+1)^c+1\le (C+1)(n+1)^c$.
\end{lemma}

\begin{definition}[Bounded search for the input/certificate split]
\lean{Complexity.certificateSplit}
\label{Complexity.certificateSplit}
\leanok
For $C, c, m \in \mathbb{N}$, the partial value $\mathrm{split}_{C,c}(m)$ is the
least $n \in \{0, 1, \dots, m\}$ satisfying $n + (C+1)(n+1)^c = m$, found by
searching $0, 1, \dots, m$ in increasing order; it is undefined when no such $n$
exists.
\end{definition}

\begin{lemma}[Specification of certificate split search]
\lean{Complexity.certificateSplit_spec}
\label{Complexity.certificateSplit_spec}
\leanok
\uses{Complexity.certificateSplit, Complexity.certificateTotal_strictMono}
\difficulty{4}
Let $C,c,m,n$ be natural numbers. The bounded certificate split search at length $m$
returns $n$ if and only if $n+(C+1)(n+1)^c=m$.
\end{lemma}

\begin{lemma}[No split for zero input]
\lean{Complexity.certificateSplit_zero}
\label{Complexity.certificateSplit_zero}
\leanok
\uses{Complexity.certificateSplit, Complexity.certificateSplit_spec, Complexity.certificate_room}
\difficulty{3}
For all $C,c\in\mathbb{N}$, the value $\mathrm{split}_{C,c}(0)$ is undefined;
equivalently, there is no $n\in\{0\}$ such that $n+(C+1)(n+1)^c=0$.
\end{lemma}

\begin{definition}[Paired exact-length verifier]
\lean{Complexity.pairedVerifier}
\label{Complexity.pairedVerifier}
\leanok
\uses{Turing.pairDecode}
For $C, c \in \mathbb{N}$ and a language $V \subseteq \{0,1\}^{*}$, the language
$V^{\mathrm{pair}}_{C,c}$ consists of the binary strings $y$ that the aligned
pair parser decodes as a pair $(x, u)$ --- that is, $y = \langle x, u\rangle$ is
a well-formed self-delimiting pair encoding, consisting of $x$ with every bit
doubled, then the separator $01$, then $u$ verbatim --- such that
$|u| = C(|x|+1)^c$ and the concatenation $xu$ lies in $V$.
\end{definition}

\begin{lemma}[Membership of an Encoded Pair]
\lean{Complexity.pairedVerifier_pair}
\label{Complexity.pairedVerifier_pair}
\leanok
\uses{Complexity.pairedVerifier, Turing.pairDecode, Turing.pairDecode_pairEncode, Turing.pairEncode}
\difficulty{2}
Let $C,c\in\mathbb{N}$, let $V$ be a language of binary strings, and let $x,u$ be binary
strings. The encoding of $(x,u)$ belongs to the paired verifier language exactly when
$|u|=C(|x|+1)^c$ and the concatenation $xu$ belongs to $V$.
\end{lemma}

\begin{lemma}[Malformed pairs are rejected]
\lean{Complexity.pairedVerifier_malformed}
\label{Complexity.pairedVerifier_malformed}
\leanok
\uses{Complexity.pairedVerifier, Turing.pairDecode}
\difficulty{2}
Let $C,c\in\mathbb{N}$, let $V$ be a language of finite bit strings, and let $y$ be a
finite bit string. If the aligned pair parser is undefined on $y$, then $y$ does not
belong to the paired verifier language defined by $C,c,$ and $V$.
\end{lemma}

\begin{definition}[Padded exact-length verifier]
\lean{Complexity.paddedVerifier}
\label{Complexity.paddedVerifier}
\leanok
\uses{Complexity.certificateSplit, Complexity.stripCertificate, Turing.pairEncode}
For $C, c \in \mathbb{N}$ and a language $V \subseteq \{0,1\}^{*}$, the language
$V^{\mathrm{pad}}_{C,c}$ consists of the binary strings $y$ for which there is a
least $n \le |y|$ with $n + (C+1)(n+1)^c = |y|$, and, writing $y = xv$ with
$|x| = n$, the suffix $v$ contains a $1$, the prefix $u$ of $v$ preceding its last
$1$ satisfies $|u| \le C(n+1)^c$, and $\langle x, u\rangle \in V$. Here
$\langle x, u\rangle$ is the self-delimiting pair encoding: $x$ with every bit
doubled, then the separator $01$, then $u$ verbatim.
\end{definition}

\begin{lemma}[Inputs without a valid split are rejected]
\lean{Complexity.paddedVerifier_no_split}
\label{Complexity.paddedVerifier_no_split}
\leanok
\uses{Complexity.certificateSplit, Complexity.paddedVerifier}
\difficulty{2}
Let $C,c\in\mathbb{N}$, let $V$ be a language of binary strings, and let $y$ be a binary
string. If no $n\le |y|$ satisfies $n+(C+1)(n+1)^c=|y|$, then $y$ does not belong to the
padded verifier for $C,c,V$.
\end{lemma}

\begin{lemma}[Exact-boundary padded verifier characterization]
\lean{Complexity.paddedVerifier_append}
\label{Complexity.paddedVerifier_append}
\leanok
\uses{Complexity.certificateSplit, Complexity.certificateSplit_spec, Complexity.paddedVerifier, Complexity.stripCertificate, Turing.pairEncode}
\difficulty{3}
Let $C,c\in\mathbb{N}$, let $V$ be a language of binary strings, and let $x,v$ be binary
strings with $|v|=(C+1)(|x|+1)^c$. Then $xv$ belongs to the padded verifier for $C,c,V$
if and only if there is a binary string $u$ such that stripping the certificate from $v$
yields $u$, $|u|\le C(|x|+1)^c$, and the self-delimiting pair encoding of $x$ and $u$
belongs to $V$.
\end{lemma}

\begin{lemma}[All-zero padding is rejected]
\lean{Complexity.paddedVerifier_no_marker}
\label{Complexity.paddedVerifier_no_marker}
\leanok
\uses{Complexity.paddedVerifier, Complexity.paddedVerifier_append, Complexity.stripCertificate_false}
\difficulty{2}
For every $C,c\in\mathbb{N}$, every language $V$ of binary strings, and every binary
string $x$, the string formed by appending $(C+1)(|x|+1)^c$ zeros to $x$ does not belong
to the padded verifier for $C,c,V$.
\end{lemma}

\begin{lemma}[Rejection of Overlong Witnesses]
\lean{Complexity.paddedVerifier_too_long}
\label{Complexity.paddedVerifier_too_long}
\leanok
\uses{Complexity.paddedVerifier, Complexity.paddedVerifier_append, Complexity.stripCertificate_pad}
\difficulty{3}
Let $C,c,k\in\mathbb{N}$, let $V$ be a language of binary strings, and let $x,u$ be
binary strings. If $|u\,1\,0^k|=(C+1)(|x|+1)^c$ and $C(|x|+1)^c<|u|$, then
$x\,u\,1\,0^k$ does not belong to the padded verifier language $V^{\mathrm{pad}}_{C,c}$.
\end{lemma}

\begin{lemma}[Padded and Bounded Witnesses]
\lean{Complexity.paddedVerifier_witness}
\label{Complexity.paddedVerifier_witness}
\leanok
\uses{Complexity.certificate_room, Complexity.paddedVerifier, Complexity.paddedVerifier_append, Complexity.stripCertificate_pad, Turing.pairEncode}
\difficulty{4}
Let $C,c\in\mathbb{N}$, let $V$ be a language of binary strings, and let $x$ be a binary
string. There exists a binary string $v$ of length $(C+1)(|x|+1)^c$ such that $xv\in
V^{\mathrm{pad}}_{C,c}$ if and only if there exists a binary string $u$ with $|u|\le
C(|x|+1)^c$ and $\langle x,u\rangle\in V$, where $\langle x,u\rangle$ is obtained by
doubling each bit of $x$, appending $01$, and then appending $u$.
\end{lemma}

\begin{definition}[Guarded pair-to-concatenation map]
\lean{Complexity.verifier_concat}
\label{Complexity.verifier_concat}
\leanok
\uses{Turing.pairDecode}
The function on binary strings that sends $z$ to the concatenation $ab$ when the
aligned pair parser decodes $z$ as a pair $(a, b)$ --- that is, when
$z = \langle a, b\rangle$ is a well-formed self-delimiting pair encoding ($a$ with
every bit doubled, then $01$, then $b$) --- and sends every other string to the
empty string.
\end{definition}

\begin{lemma}[Equality of two split searches]
\lean{Complexity.verifier_split_bridge}
\label{Complexity.verifier_split_bridge}
\leanok
\uses{Complexity.certificateSplit, Turing.solveSplit}
\difficulty{1}
For all $C,c\in\mathbb{N}$, the partial functions on natural numbers given by the least
$i\le m$ satisfying $i+(C+1)(i+1)^c=m$, when such an $i$ exists, and undefined
otherwise, are equal.
\end{lemma}

\begin{lemma}[Equality of two stripping procedures]
\lean{Complexity.verifier_strip_bridge}
\label{Complexity.verifier_strip_bridge}
\leanok
\uses{Complexity.stripCertificate, Complexity.stripCertificate_spec, Turing.splitAtLastTrue}
\difficulty{4}
The partial function that returns the prefix before the last $1$ in a binary string, and
is undefined when there is no $1$, equals the partial function that recursively strips a
marked certificate by the same rule.
\end{lemma}

\begin{definition}[Original-bound test on encoded pairs]
\lean{Complexity.verifier_bound}
\label{Complexity.verifier_bound}
\leanok
\uses{Turing.pairDecode}
For $C, c \in \mathbb{N}$, the Boolean function $\beta_{C,c}$ on binary strings
that, on a string $z$ which the aligned pair parser decodes as a pair $(a, b)$ ---
that is, $z = \langle a, b\rangle$ is a well-formed self-delimiting pair encoding
($a$ with every bit doubled, then $01$, then $b$) --- returns true exactly when
$|b| \le C(|a|+1)^c$, and returns false on every string that is not a
well-formed pair encoding.
\end{definition}

\begin{lemma}[Polynomial bound test is polynomial-time computable]
\lean{Complexity.verifier_poly_bound}
\label{Complexity.verifier_poly_bound}
\leanok
\uses{Complexity.PolyTimeComputable, Complexity.verifier_bound, Turing.FinTM.computesFunInTime_pairLenCheck}
\difficulty{1}
For any natural numbers $C$ and $c$, the function that sends a bit string $z$ to the
one-bit string indicating whether its decoded pair $(a,b)$ satisfies $|b|\le C(|a|+1)^c$
is polynomial-time computable; strings that do not decode as pairs are assigned
indicator value $0$.
\end{lemma}

\begin{lemma}[Polynomial-time language indicators]
\lean{Complexity.verifier_poly_indicator}
\label{Complexity.verifier_poly_indicator}
\leanok
\uses{Complexity.P, Complexity.PolyTimeComputable, Complexity.mem_P_iff, Complexity.polyTimeComputable_const, Complexity.polyTimeComputable_ite, Turing.MultiTapeTM.indicator}
\difficulty{3}
Let $V$ be a language over the binary alphabet and suppose $V$ belongs to $mathsf{P}$.
Then the function sending each binary string $x$ to the one-bit string containing $1$ if
$x\in V$ and $0$ otherwise is polynomial-time computable.
\end{lemma}

\begin{lemma}[Polynomial-time indicators define a language in P]
\lean{Complexity.verifier_mem_P}
\label{Complexity.verifier_mem_P}
\leanok
\uses{Complexity.P, Complexity.PolyTimeComputable, Complexity.mem_P_iff, Turing.MultiTapeTM.indicator}
\difficulty{1}
Let $V$ be a language of bit strings. If the function that sends each bit string $x$ to
the one-bit string indicating whether $x\in V$ is polynomial-time computable, then $V$
belongs to $\mathsf{P}$.
\end{lemma}

\begin{lemma}[Polynomial-time reverse length comparison]
\lean{Complexity.verifier_poly_reverseBound}
\label{Complexity.verifier_poly_reverseBound}
\leanok
\uses{Complexity.PolyTimeComputable, Complexity.PolyTimeComputable.comp, Complexity.PolyTimeComputable.pairEncode, Complexity.pairFstD, Complexity.pairSndD, Complexity.polyTimeComputable_of_linear, Complexity.polyTimeComputable_pairFstD, Complexity.polyTimeComputable_pairSndD, Complexity.verifier_bound, Complexity.verifier_poly_bound, Turing.FinTM.computesFunInTime_polyUnary, Turing.FinTM.computesFunInTime_prepend, Turing.pairDecode_pairEncode}
\difficulty{5}
For all $C,c\in\mathbb{N}$, the function on binary strings that sends $z$ to the one-bit
string containing $1$ exactly when $C(|\pi_1(z)|+1)^c\le |\pi_2(z)|$, and containing $0$
otherwise, is polynomial-time computable. Here each projection returns the corresponding
component when $z$ is a valid encoded pair and the empty string otherwise.
\end{lemma}

\begin{lemma}[Polynomial-time closure of the paired verifier]
\lean{Complexity.pairedVerifier_mem_P}
\label{Complexity.pairedVerifier_mem_P}
\leanok
\uses{Complexity.P, Complexity.PolyTimeComputable, Complexity.PolyTimeComputable.comp, Complexity.pairFstD, Complexity.pairSndD, Complexity.pairedVerifier, Complexity.polyTimeComputable_const, Complexity.polyTimeComputable_ite, Complexity.polyTimeComputable_of_linear, Complexity.verifier_bound, Complexity.verifier_concat, Complexity.verifier_mem_P, Complexity.verifier_poly_bound, Complexity.verifier_poly_indicator, Complexity.verifier_poly_reverseBound, Turing.FinTM.computesFunInTime_pairConcat, Turing.FinTM.computesFunInTime_pairValid, Turing.MultiTapeTM.indicator, Turing.pairDecode}
\difficulty{5}
Let $C,c\in\mathbb{N}$ and let $V$ be a language of finite bit strings. If $V$ is
decidable in polynomial time, then the paired verifier language defined by $C,c,$ and
$V$ is also decidable in polynomial time.
\end{lemma}

\begin{definition}[Split of a padded input into a pair]
\lean{Complexity.verifier_split}
\label{Complexity.verifier_split}
\leanok
\uses{Turing.pairEncode, Turing.solveSplit}
For $C, c \in \mathbb{N}$, the function on binary strings that, given $y$, finds
the least $n \le |y|$ with $n + (C+1)(n+1)^c = |y|$; if such an $n$ exists it
outputs the self-delimiting pair encoding $\langle x, v\rangle$ ($x$ with every
bit doubled, then $01$, then $v$), where $x$ is the first $n$ bits of $y$ and $v$
the remaining bits, and otherwise it outputs the empty string.
\end{definition}

\begin{definition}[Stripping the payload of an encoded pair]
\lean{Complexity.verifier_strip}
\label{Complexity.verifier_strip}
\leanok
\uses{Turing.pairDecode, Turing.pairEncode, Turing.splitAtLastTrue}
The function on binary strings that, on a well-formed self-delimiting pair
encoding $z = \langle a, v\rangle$ ($a$ with every bit doubled, then $01$, then
$v$) in which $v$ contains a $1$, outputs $\langle a, u\rangle$, where $u$ is the
prefix of $v$ preceding its last $1$; on every other string (a malformed pair, or
a pair whose second component has no $1$) it outputs the empty string.
\end{definition}

\begin{lemma}[Polynomial-time closure of padded verification]
\lean{Complexity.paddedVerifier_mem_P}
\label{Complexity.paddedVerifier_mem_P}
\leanok
\uses{Complexity.P, Complexity.PolyTimeComputable, Complexity.PolyTimeComputable.comp, Complexity.certificateSplit, Complexity.certificateSplit_spec, Complexity.paddedVerifier, Complexity.polyTimeComputable_const, Complexity.polyTimeComputable_ite, Complexity.polyTimeComputable_of_linear, Complexity.stripCertificate, Complexity.verifier_bound, Complexity.verifier_mem_P, Complexity.verifier_poly_bound, Complexity.verifier_poly_indicator, Complexity.verifier_split, Complexity.verifier_split_bridge, Complexity.verifier_strip, Complexity.verifier_strip_bridge, Turing.FinTM.computesFunInTime_pairValid, Turing.FinTM.computesFunInTime_splitSolve, Turing.FinTM.computesFunInTime_stripLast, Turing.MultiTapeTM.indicator, Turing.pairDecode, Turing.pairDecode_pairEncode}
\difficulty{6}
Let $C,c\in\mathbb{N}$ and let $V$ be a language of binary strings in polynomial time.
The padded verifier for $C,c,V$ is also in polynomial time.
\end{lemma}

\begin{theorem}[NP certificates of bounded length]
\lean{Complexity.mem_NP_iff_exists_length_le}
\label{Complexity.mem_NP_iff_exists_length_le}
\leanok
\uses{Complexity.NP, Complexity.P, Complexity.paddedVerifier, Complexity.paddedVerifier_mem_P, Complexity.paddedVerifier_witness, Complexity.pairedVerifier, Complexity.pairedVerifier_mem_P, Complexity.pairedVerifier_pair, Turing.pairEncode}
\difficulty{4}
A language $L$ over the binary alphabet belongs to $\mathsf{NP}$ if and only if there
are natural numbers $C,c$ and a language $V\in\mathsf{P}$ such that, for every binary
string $x$, $x\in L$ exactly when there is a binary string $u$ with $|u|\le C(|x|+1)^c$
for which $\langle x,u\rangle\in V$. Here $\langle x,u\rangle$ is formed by doubling
each bit of $x$, appending the separator $01$, and then appending $u$.
\end{theorem}
